
On
considère alors le morphisme canonique (2.7.2) :
\[
  \underline{N}_1^{\bullet}=\operatorname{Comp}_{d_0}(\underline{M}^{\bullet})\otimes_{\underline{O}_S}\underline{B}(-d_0)
  \longrightarrow \underline{M}^{\bullet};
\]
soit $\underline{P}_1^{\bullet}$ le cône de ce morphisme. C'est encore un complexe $\mathcal C_0$-parfait : en effet, pour tout $n$, $\operatorname{Comp}_n(\underline{M}^{\bullet})$ est un complexe parfait sur $S$, car il en est ainsi quand $\underline{M}^{\bullet}$ est borné à termes localement libres de type fini sur $\underline{B}$; par suite, le complexe $\operatorname{Comp}_n(\underline{M}^{\bullet})\otimes_{\underline{O}_S}\underline{B}(-d_0)$ est $\mathcal C_0$-parfait, et il en est donc de même pour le cône (I 4.10).

Par ailleurs, on peut écrire sur tout ouvert $U$ du recouvrement le diagramme commutatif :
\[
\begin{tikzcd}[column sep=large,row sep=large]
\operatorname{Comp}_{d_0}(\underline{M}^{\bullet})\otimes_{\underline{O}_S}\underline{B}(-d_0) \arrow[r] \arrow[d,dash,"\wr"] & \underline{M}^{\bullet} \arrow[d,dash,"\wr"]\\
\operatorname{Comp}_{d_0}(\underline{M}_U'{}^{\bullet})\otimes_{\underline{O}_S}\underline{B}(-d_0) \arrow[r] & \underline{M}_U'{}^{\bullet},
\end{tikzcd}
\]
qui induit un isomorphisme entre les cônes; or l'homomorphisme du bas est un isomorphisme en degré $d_0$; donc $\underline{P}_1^{\bullet}$ est isomorphe sur $U$ à un complexe borné à termes localement libres de type fini sur $\underline{B}$, dont les composantes $\operatorname{Comp}_d$ sont nulles pour $d<d_0+1$. Par suite, $\underline{P}_1^{\bullet}$ vérifie les mêmes propriétés que $\underline{M}^{\bullet}$, $d_0$ étant remplacé par $d_0+1$. On peut donc itérer le processus, et il suffit de voir qu'au bout d'un nombre fini d'opérations, le complexe $\underline{P}_k^{\bullet}$ obtenu est acyclique. Là encore la question est locale, si bien qu'il suffit de le voir pour $\underline{M}_U'{}^{\bullet}$. Or pour chaque terme $\underline{M}_U'{}^i$, les morphismes considérés plus haut sont ceux qui définissent la filtration de la proposition 1.7; cette filtration est finie, donc en un nombre fini d'opérations, on arrive à partir de $\underline{M}_U'{}^{\bullet}$ au complexe nul, c'est-à-dire que le complexe $\underline{P}_k^{\bullet}$ correspondant est acyclique.

